% ACM Conference/Journal Template
% Usage: Replace placeholders with your content
% Requires: acmart.cls (standard in TeX Live / MiKTeX)

\documentclass[sigconf,review,anonymous]{acmart}
% Options: sigconf (conference), acmtog (journal), acmlarge, acmsmall
% Add 'anonymous' for double-blind, 'review' for review mode

\usepackage{booktabs}
\usepackage{algorithm}
\usepackage{algorithmic}
\usepackage{multirow}

% Remove ACM-specific metadata for draft
\setcopyright{none}
\settopmatter{printacmref=false}
\renewcommand\footnotetextcopyrightpermission[1]{}
\pagestyle{plain}

\begin{document}

\title{TODO: Paper Title}

\author{TODO: First Author}
\authornote{TODO: Note if needed}
\email{todo@university.edu}
\orcid{TODO: 0000-0000-0000-0000}
\affiliation{%
  \institution{TODO: University}
  \streetaddress{TODO: Address}
  \city{TODO: City}
  \country{TODO: Country}
}

\author{TODO: Second Author}
\email{todo@university.edu}
\affiliation{%
  \institution{TODO: University}
  \city{TODO: City}
  \country{TODO: Country}
}

\begin{abstract}
TODO: Write abstract (150-250 words).
Structure: Problem -- Gap -- Method -- Key Result -- Implication.
\end{abstract}

% ACM Computing Classification System
\begin{CCSXML}
<ccs2012>
  <concept>
    <concept_id>TODO: CCS concept ID</concept_id>
    <concept_desc>TODO: CCS concept description</concept_desc>
    <concept_significance>500</concept_significance>
  </concept>
</ccs2012>
\end{CCSXML}

\ccsdesc[500]{TODO: Computing methodologies~Machine learning}

\keywords{TODO: keyword1, keyword2, keyword3}

\maketitle

% ============================================================
\section{Introduction}
\label{sec:introduction}

TODO: Motivation, gap, approach, main findings, contributions.

Our contributions are:
\begin{itemize}
  \item TODO: Contribution 1 (specific, evidence-backed)
  \item TODO: Contribution 2
  \item TODO: Contribution 3
\end{itemize}

% ============================================================
\section{Related Work}
\label{sec:related_work}

\subsection{TODO: Cluster 1}
TODO: Thematic group with contrast to your work.

\subsection{TODO: Cluster 2}
TODO: Second thematic group.

% ============================================================
\section{Method}
\label{sec:method}

\subsection{Problem Setup}
TODO: Formal problem statement, notation.

\subsection{Proposed Approach}
TODO: Detailed method description.

% ============================================================
\section{Experiments}
\label{sec:experiments}

\subsection{Setup}
TODO: Datasets, baselines, metrics, implementation details.

\subsection{Results}
TODO: Main comparison results.

\begin{table}[htbp]
\caption{TODO: Main Results}
\label{tab:main_results}
\begin{tabular}{lccc}
\toprule
Method & Metric 1 & Metric 2 & Metric 3 \\
\midrule
Baseline 1 & TODO & TODO & TODO \\
Baseline 2 & TODO & TODO & TODO \\
\textbf{Ours} & \textbf{TODO} & \textbf{TODO} & \textbf{TODO} \\
\bottomrule
\end{tabular}
\end{table}

\subsection{Ablation Study}
TODO: Component analysis.

% ============================================================
\section{Discussion and Limitations}
\label{sec:discussion}
TODO: Interpret results, limitations, future work.

% ============================================================
\section{Conclusion}
\label{sec:conclusion}
TODO: Summary of contributions and key findings.

\begin{acks}
TODO: Acknowledgments, funding, compute resources.
\end{acks}

\bibliographystyle{ACM-Reference-Format}
\bibliography{references}

\end{document}
